\section{A concentrating solution on the compact manifold}
\label{glue:section}
\label{sec:compact}

We now correct the flat solution on the auxiliary sphere.  All functions and
operators in this Section are restricted to the
$O(m)\times O(6)$ invariant subspace, where $m=24$ and
$\Sigma=S^{29}\subset\R^{24}\times\R^6$.  We use the spaces of
Definition~\ref{ana:spaces}, on $\Sigma\times[-T,0)$, with
\[
 r=\sqrt{x},\qquad \rho=\sqrt{-t},\qquad
 \sigma=(r^2+\rho)^{1/2},\qquad p_0=-2+2\delta.
\]
When necessary, a subscript specifies the H\"older exponent or the scale:
for example, $X_\alpha(s)$ has weight $p_0$ and scale $s$, and
$Y_\alpha(s)$ has weight $p_0-4$.  Without a specified scale, it is
$\sigma$.  The space $X'_\alpha$ includes the parabolic
$C^\alpha$ norm of every spatial derivative of order at most four; only
the time derivative is omitted.

\begin{theorem}[Compact realization]\label{glue:solution}
Let $G$ be the profile of Theorem~\ref{prof:profile}, put
\[
 h(y)=\frac{G(y)-1}{y},\qquad
 F(t,x)=\rho^{-1}h(x/\rho),
\]
and fix $0<\alpha<\alpha'<1$.  Choose $\delta\in(0,1)$ and the ancient
right inverse supplied by Theorem~\ref{flat:inverse}.  For every
sufficiently small $\eps>0$, there are $T>0$ and a real invariant
function $V\in X_\alpha$ on $\Sigma\times[-T,0)$ such that
$\norm{V}_{X'_\alpha}\le\eps$ and $H=F+V$ satisfies
\begin{equation}\label{glue:equation}
 H_t+(1+v_*H)^2\A H=0,\qquad v_*=x(1-x).
\end{equation}
The function $H$ is smooth on $\Sigma\times(-T,0)$, and there are
constants $0<c_1<c_2<\infty$, independent of $t$, such that
\begin{equation}\label{glue:positive}
 c_1\le 1+v_*H\le c_2.
\end{equation}
In particular, $H(t,\cdot)$ is smooth as a function of $x\in[0,1]$, and
\begin{equation}\label{glue:center}
 V(t,0)=O(\rho^{-1+\delta})\qquad(t\uparrow0).
\end{equation}
\end{theorem}

We first establish the estimates and the right inverse needed in the
proof.  All constants below may depend on the fixed profile, the
exponents, the sphere, and the fixed box radius in
Definition~\ref{ana:spaces}.  They do not depend on the localization
radius $\beta$ or on $T$, subject to the restrictions explicitly stated.
Changing that fixed box radius to another sufficiently small fixed
radius changes the norms by uniform constants, by the covering property
in Definition~\ref{ana:spaces}.

\subsection{Background estimates and the compact error}

Write
\begin{align*}
 g_c&=1+v_*F=x+(1-x)G(x/\rho),&
 A_c&=g_c^2,& B_c&=2v_*g_c\A F,\\
 g_f&=1+xF=G(x/\rho),&
 A_f&=g_f^2,& B_f&=2xg_f\wh\A F.
\end{align*}
Here $B_f$ is in the original flat variables, before the similarity
change of variables in Section~\ref{sec:flat}.  Set
\begin{equation}\label{glue:linearizations}
 \Pcal=\partial_t+A_c\A+B_c,\qquad
 \wh\Pcal=\partial_t+A_f\wh\A+B_f,\qquad
 E_c=F_t+A_c\A F.
\end{equation}
The profile equation gives $F_t+A_f\wh\A F=0$.  Moreover, if
\[
 \mu=\min(1,\inf G)>0,\qquad M=\max(1,\sup G),
\]
then
\begin{equation}\label{glue:background-positive}
 \mu\le g_c\le M.
\end{equation}

We make the meaning of all the scaled estimates explicit.  In a box
with center $(z_0,t_0)$ and center scale $b=\sigma(z_0,t_0)$,
$O_{\mathrm{sc}}(b^q)$ means that, for every pair of nonnegative integers
$j,k$, the spatial derivatives of order $j$ and time derivatives of
order $k$ are bounded by $C_{j,k}b^{q-j-4k}$ on that box.  The scale $b$
is held fixed when differentiating.  These bounds also give the
corresponding scaled parabolic H\"older bounds of any exponent less
than one, on slightly smaller boxes.  The notation will only be used
for smooth background quantities, not for arbitrary members of $X$
or $Y$.

\begin{lemma}[Scaled background and mismatch]\label{glue:background}
For $0<T\le1$, on the compact domain one has
\begin{equation}\label{glue:background-bounds}
 F=O_{\mathrm{sc}}(b^{-2}),\qquad
 \A F=O_{\mathrm{sc}}(b^{-6}),\qquad
 A_c=O_{\mathrm{sc}}(1),\qquad
 B_c=O_{\mathrm{sc}}(b^{-4}),\qquad
 v_*=O_{\mathrm{sc}}(b^2).
\end{equation}
The analogous flat bounds hold in a fixed neighborhood of $r=0$.
In the common coordinates there, write
\[
 \Pcal-\wh\Pcal=\sum_{|\nu|\le4}d_\nu(z,t)\partial_z^\nu.
\]
For each $j,k\ge0$, the rescaled coefficients satisfy
\begin{equation}\label{glue:mismatch-coeff}
 b^{4-|\nu|+j+4k}
 \abs{\partial_z^j\partial_t^k d_\nu}\le C_{j,k}b^2.
\end{equation}
Finally, near $r=0$,
\begin{equation}\label{glue:error-order}
 E_c=O_{\mathrm{sc}}(b^{-4}).
\end{equation}
Away from a fixed neighborhood of $r=0$, all these background
quantities extend smoothly to $t=0$ with bounds of every finite order.
\end{lemma}

\begin{proof}
In the region $r^2\le C\rho$, the formula for $F$, smoothness of $h$ at
zero, and the change $z=\sqrt\rho\,\zeta$ give
\[
 \abs{\partial_z^j\partial_t^k F}
 \le C_{j,k}\rho^{-1-j/2-2k}
 \le C'_{j,k}\sigma^{-2-j-4k}.
\]
In the region $r^2\ge\rho$, use the differentiable expansion from
Theorem~\ref{prof:profile}.  For arbitrary $N$ it has the form
\begin{equation}\label{glue:outer-expansion}
 F(t,r^2)=\frac{c-1}{r^2}
       +\sum_{\ell=1}^{N}a_\ell\frac{(-t)^\ell}{r^{4\ell+2}}
       +\text{remainder}.
\end{equation}
The all-order remainder estimates allow $N$ to be chosen larger than
any prescribed number of differentiations.  Differentiating the
displayed terms and the remainder gives
$\abs{\partial_z^j\partial_t^kF}\le C_{j,k}r^{-2-j-4k}$.
This proves the first bound and also proves the asserted smooth
extension away from the core.  In particular it supplies time
derivative bounds up to the open terminal face, which do not follow
from spatial decay alone.

Near the core, write a point of the sphere as
$(z,\sqrt{1-|z|^2}\,\theta)$, $\theta\in S^5$.  Invariant functions
are independent of $\theta$, and on these functions
\begin{equation}\label{glue:laplace-coordinate}
 \Delta_\Sigma
 =\Delta_z-\sum_{i,j}z_i z_j\partial_{ij}
             -29\sum_i z_i\partial_i.
\end{equation}
The product coordinates and ordinary sphere coordinates have
uniformly controlled transition maps on a fixed smaller
neighborhood.  Thus their scaled norms are uniformly equivalent for
the functions considered here.  Lemma~\ref{geo:operators} gives
\[
 \A=\frac1{16}\Delta_\Sigma^2
       -\frac{37}{4}\Delta_\Sigma+312,
 \qquad \wh\A=\frac1{16}\Delta_z^2.
\]
In $\A-\wh\A$, the coefficients of orders four, three, and two have
sizes $O(r^2)$, $O(r)$, and $O(1)$, respectively; those of orders one
and zero are smooth and bounded.  After multiplication by $b^4$ and
rescaling space by $b$, every coefficient is $O_{\mathrm{sc}}(b^2)$.
For example the first three sizes become $O(r^2)$, $O(br)$, and
$O(b^2)$, and $r\le b$.  Differentiating these polynomial coefficients
preserves the stated scaled bound.  The bounds for $\A F$ follow;
away from this coordinate neighborhood they follow from the smooth
extension and the positive lower bound for $\sigma$.

The function $v_*=x-x^2$ has size $O_{\mathrm{sc}}(b^2)$, and
$v_*-x=-x^2=O_{\mathrm{sc}}(b^4)$.  Hence
\[
 g_c-g_f=-x^2F=O_{\mathrm{sc}}(b^2),\qquad
 A_c-A_f=O_{\mathrm{sc}}(b^2).
\]
Products give the remaining bounds in
\eqref{glue:background-bounds}.  For the zeroth-order mismatch, use
\[
 B_c-B_f
 =2(v_*g_c-xg_f)\wh\A F
        +2v_*g_c(\A-\wh\A)F.
\]
The two terms have size $O_{\mathrm{sc}}(b^{-2})$, since
$v_*g_c-xg_f=O_{\mathrm{sc}}(b^4)$ and
$(\A-\wh\A)F=O_{\mathrm{sc}}(b^{-4})$.
Together with the principal multiplier difference, this proves
\eqref{glue:mismatch-coeff} in every derivative order.  Finally,
\[
 E_c=(A_c-A_f)\wh\A F+A_c(\A-\wh\A)F
      =O_{\mathrm{sc}}(b^{-4}),
\]
using the exact flat equation, as required.
\end{proof}

One consequence that we will use repeatedly is that the local
$Y_\vartheta$ size of $E_c$, for either $\vartheta=\alpha$ or
$\vartheta=\alpha'$, is at most
\begin{equation}\label{glue:error-weight}
 C b^{6-2\delta}b^{-4}=C b^{2-2\delta}
\end{equation}
near the core.  The same statement holds with any fixed finite number
of extra scaled derivatives.  Globally $\norm{E_c}_{Y_\vartheta}\le C$.
If it is multiplied by a smooth cutoff supported in $r\le C_0\beta$
and varying only where $r$ is comparable to $\beta$, then, whenever
$T\le\beta^4$,
\begin{equation}\label{glue:core-error-small}
 \norm{\text{cutoff}\,E_c}_{Y_\vartheta}
 \le C\beta^{2-2\delta}.
\end{equation}
Indeed every box meeting this support has $b\le C\beta$; cutoff
derivatives have bounded scaled size.  The strict inequality
$\delta<1$ is used here.

\subsection{Localization of the ancient inverse}

Multiplying a flat correction by a spatial cutoff produces a commutator
on the transition annulus.  This error is bounded but need not be small.
We will include its negative in the forcing for the outer forward
problem, so it cancels in the combined equation.  The errors that remain
are the geometric mismatch, made small by shrinking the core, and the
outer cutoff commutator, made small by shortening the time interval.
The nested cutoffs below ensure that these cancellations are exact.

Choose once and for all a smooth function
$\chi:[0,\infty)\to[0,1]$, equal to one on $[0,1]$ and zero on
$[2,\infty)$.  All its compositions below are smooth at $r=0$
because they are constant there.  Take $\beta>0$ so small that
$r\le5\beta$ lies in the common coordinate neighborhood above, and
assume $T\le\beta^4$.  Define
\begin{equation}\label{glue:cutoffs}
 \chi_i=\chi(r/\beta),\qquad
 \chi_b=\chi(r/(2\beta)),\qquad
 \chi_o=1-\chi(2r/\beta),\qquad
 \kappa=1-\chi(4r/\beta).
\end{equation}
Thus the inner input cutoff $\chi_i$ is supported in $r\le2\beta$;
the inner output cutoff $\chi_b$ is one there and supported in
$r\le4\beta$.  The outer output cutoff $\chi_o$ is zero for
$r\le\beta/2$ and one for $r\ge\beta$.  The coefficient cutoff
$\kappa$ is already one for $r\ge\beta/2$ and is zero for
$r\le\beta/4$.

Figure~\ref{fig:cutoffs} records these nested regions. In particular,
the inner output cutoff is one on the support of the inner input,
and the modified outer coefficients agree with the compact coefficients
where the outer output cutoff can be nonzero. The cutoff commutators are supported on the corresponding transition
annuli.

\begin{figure}[htbp]
\centering
\begin{tikzpicture}[x=.96cm,y=.6cm,font=\small]
\definecolor{cutone}{RGB}{40,83,111}
\tikzset{oneregion/.style={fill=cutone,draw=cutone},
 transition/.style={pattern=north east lines,pattern color=black!55,draw=black!55}}
\foreach \x/\lab in {0/0,1/{\beta/4},2/{\beta/2},3/{\beta},4.5/{2\beta},6/{4\beta}} {
  \draw[black!18] (\x,.45)--(\x,4.6);
  \node[above] at (\x,4.6) {$\lab$};
}
\draw[->] (0,4.6)--(6.8,4.6) node[right] {$r$};
\foreach \yy/\lab in {4/{\chi_i},3/{\chi_b},2/{\chi_o},1/{\kappa}} {
 \draw[black!30] (0,\yy-.18) rectangle (6.6,\yy+.18);
 \node[left] at (-.25,\yy) {$\lab$};
}
\filldraw[oneregion] (0,3.82) rectangle (3,4.18);
\filldraw[transition] (3,3.82) rectangle (4.5,4.18);
\filldraw[oneregion] (0,2.82) rectangle (4.5,3.18);
\filldraw[transition] (4.5,2.82) rectangle (6,3.18);
\filldraw[transition] (2,1.82) rectangle (3,2.18);
\filldraw[oneregion] (3,1.82) rectangle (6.6,2.18);
\filldraw[transition] (1,.82) rectangle (2,1.18);
\filldraw[oneregion] (2,.82) rectangle (6.6,1.18);
\filldraw[oneregion] (-.3,-.13) rectangle (.15,.13);
\node[right] at (.22,0) {identically $1$};
\filldraw[transition] (2.35,-.13) rectangle (2.8,.13);
\node[right] at (2.87,0) {may vary};
\draw[black!30] (4.85,-.13) rectangle (5.3,.13);
\node[right] at (5.37,0) {identically $0$};
\end{tikzpicture}

\caption{The cutoff regions in \eqref{glue:cutoffs}. Distances are
schematic. Solid bands indicate the value one, blank bands the value
zero, and hatched bands the intervals where a cutoff may vary in
$[0,1]$; no shape or monotonicity of the transition is assumed.}
\label{fig:cutoffs}
\end{figure}

\begin{lemma}[Extension of the inner source]\label{glue:extension}
There is a real linear extension $\mathcal E_i:Y_\alpha\to
Y_\alpha(\R^m\times(-\infty,0))$, uniformly bounded for
$T\le\beta^4$, which equals $\chi_i f$ in flat coordinates for
$-T\le t<0$, and is supported spatially in $r\le2\beta$.
\end{lemma}

\begin{proof}
Put $f_i=\chi_i f$, view it in the flat coordinates, and extend it by
zero beyond their neighborhood.  The spatial extension is bounded
because the cutoff vanishes before the boundary of that neighborhood.
The $Y_\alpha$ norm gives a spatial $C^\alpha$ trace $f_i(-T,z)$ at
the lower time face.  Choose a smooth $\psi:[0,\infty)\to[0,1]$ equal
to one near zero and zero for arguments at least one.  For $t<-T$ set
\begin{equation}\label{glue:extension-formula}
 (\mathcal E_i f)(t,z)
 =f_i(-T,z)\,\psi\left(\frac{-t-T}{(|z|^2+\sqrt T)^2}\right).
\end{equation}
Let $b_0=(|z|^2+\sqrt T)^{1/2}$.  Wherever this expression is
nonzero,
\[
 T\le-t\le T+b_0^4,
 \qquad b_0^2\le |z|^2+\sqrt{-t}\le(1+\sqrt2)b_0^2.
\]
The ancient scale is therefore comparable to $b_0$.  The argument of
$\psi$ has bounded derivatives in spatial units $b_0$ and temporal
units $b_0^4$ on its transition region.  The trace norm, H\"older
multiplication, and scale comparability prove the required estimate
below $-T$.  Across that face the values agree.  Splitting a temporal
increment at $-T$ bounds its $\alpha/4$ H\"older quotient by the sum
of the two one-sided bounds; a spatial increment is handled by the
trace bound on comparable boxes.  This proves the full parabolic
$C^\alpha$ bound.  No time derivative of the source is required.
The construction is real, linear, and radial.  Its continuation below
$-T$ is in fact supported in $-26\beta^4\le t<-T$, though a temporal
support restriction is unnecessary for the ancient inverse.
\end{proof}

Apply Theorem~\ref{flat:inverse} to obtain
\[
 W_f=\wh D\mathcal E_i f,\qquad V_i=\chi_b W_f,
\]
where $V_i$ is regarded as a function on the sphere and extended by
zero.  Invariant norm equivalence and the cutoff bounds give
$\norm{V_i}_{X_\alpha}\le C\norm{f}_{Y_\alpha}$.  Define
\begin{equation}\label{glue:inner-errors}
 R_i f=(\Pcal-\wh\Pcal)(\chi_bW_f),\qquad
 J_i f=[\wh\Pcal,\chi_b]W_f.
\end{equation}
Both expressions are taken in the common coordinates and extended by
zero.  Since $\chi_b f_i=f_i$, the exact identity is
\begin{equation}\label{glue:inner-identity}
 \Pcal V_i=f_i+R_i f+J_i f.
\end{equation}
Lemma~\ref{glue:background} and $b\le C\beta$ on the relevant boxes
give
\begin{equation}\label{glue:inner-error-bounds}
 \norm{R_i f}_{Y_\alpha}\le C\beta^2\norm{f}_{Y_\alpha},
 \qquad
 \norm{J_i f}_{Y_\alpha}\le C\norm{f}_{Y_\alpha}.
\end{equation}
Here $R_i f$ is supported in $r\le4\beta$, and $J_i f$ in
$2\beta\le r\le4\beta$.  The commutator has order at most three;
on its support $b$ is comparable to $\beta$, so every derivative of
$\chi_b$ is uniformly bounded in scale units.  The definition of
$R_i$ includes the mismatch applied to the cutoff itself, and hence
accounts for all mixed cutoff and mismatch terms.

\subsection{The outer Cauchy problem and the exact right inverse}

Set
\[
 s_\beta=\max(\beta,r),\qquad \eta=T/\beta^4,
\]
and use the nonsingular weighted spaces with scale $s_\beta$.
Equivalently one can use the smooth comparable scale
$(\beta^2+r^2)^{1/2}$; the norms and all following estimates are
uniformly equivalent.  Define a modified operator on the entire sphere
by
\begin{equation}\label{glue:modified-operator}
 \Pcal_\beta=\partial_t+A_\beta\A+B_\beta,
 \qquad
 A_\beta=\kappa A_c+(1-\kappa),\qquad B_\beta=\kappa B_c.
\end{equation}
Its principal multiplier is bounded above and below by positive
constants, and $\Pcal_\beta=\Pcal$ for $r\ge\beta/2$.
Lemma~\ref{glue:background} shows, for every $j,k$, uniform bounds for
the derivatives of $A_\beta$ and $s_\beta^4B_\beta$ in units
$(s_\beta,s_\beta^4)$, with the scale fixed at the center of each box.
On the transition annulus $r\asymp\beta$ these are exactly its scaled
background bounds and the cutoff bounds.  Inside $r\le\beta/4$ the
coefficients are $1$ and $0$; outside the transition they are the
original coefficients with $\sigma\asymp s_\beta$.  The sphere
coefficients of $\A$ also have uniform bounds in these units.  These
facts verify all coefficient and ellipticity hypotheses of
Lemma~\ref{ana:short}, uniformly for $\eta\le1$.

The outer source is
\begin{equation}\label{glue:outer-source}
 f_o=f-f_i-J_i f.
\end{equation}
It vanishes for $r<\beta$ and has
\begin{equation}\label{glue:outer-source-bound}
 \norm{f_o}_{Y_\alpha(s_\beta)}\le C\norm{f}_{Y_\alpha}.
\end{equation}
Indeed $s_\beta\asymp\sigma$ on its support and on every sufficiently
small box meeting that support.  The latter assertion follows from
the Lipschitz bound for $r$ and the box radius $a_0\sigma$; it also
handles boxes crossing a support boundary.  The same comparison holds
for outputs supported in $r\ge\beta/2$.

For each fixed $\beta,T$, solve the closed-manifold Cauchy problem
\begin{equation}\label{glue:outer-cauchy}
 \Pcal_\beta Z_o=f_o,\qquad Z_o(-T)=0.
\end{equation}
The modified coefficients extend smoothly to $t=0$ by
\eqref{glue:outer-expansion}, and $f_o$ extends H\"older continuously
there: for fixed $\beta$ its nonsingular weighted norm is an ordinary
H\"older bound on a compact cylinder.  Thus the solvability assertion
in Lemma~\ref{ana:short} applies, or one may solve on truncated
cylinders and use its uniform bounds.  Uniqueness makes $Z_o$ invariant
and the solution operator real linear.

Choose a fixed $\eta_0>0$ small enough for Lemma~\ref{ana:short}
at both exponents $\alpha$ and $\alpha'$.  That Lemma gives, for
$0<\eta\le\eta_0$,
\begin{align}
 \norm{Z_o}_{X_\alpha(s_\beta)}
   &\le C\norm{f_o}_{Y_\alpha(s_\beta)},\label{glue:outer-bound}\\
 \sup s_\beta^{-p_0}\abs{Z_o}
   &\le C\eta\norm{f_o}_{Y_\alpha(s_\beta)}.\label{glue:outer-sup}
\end{align}
The constants are uniform in $\beta,T$ by the coefficient and scale
bounds verified above.  The small parameter in Lemma~\ref{ana:short}
is precisely $\eta=T/\beta^4$, since $\inf s_\beta=\beta$.
That lemma applies to zero initial data even when $f_o(-T)$ is nonzero.

Put $V_o=\chi_o Z_o$.  On the support of this product the coefficients
have already become those of $\Pcal$, and $\chi_o f_o=f_o$.
Consequently
\begin{equation}\label{glue:outer-identity}
 \Pcal V_o=f_o+J_o f,\qquad
 J_o f=[\Pcal_\beta,\chi_o]Z_o.
\end{equation}
This commutator is supported in $\beta/2\le r\le\beta$ and involves
spatial derivatives of $Z_o$ of order at most three.  For these
derivatives, Lemma~\ref{ana:interpolation} and
\eqref{glue:outer-bound}--\eqref{glue:outer-sup} give the scaled
parabolic $C^\alpha$ bound
$C\eta^{(1-\alpha)/4}\norm{f_o}_{Y_\alpha(s_\beta)}$.
Since $s_\beta\asymp\sigma\asymp\beta$ on the commutator annulus,
\begin{equation}\label{glue:outer-error-small}
 \norm{J_o f}_{Y_\alpha}
 \le C\eta^{(1-\alpha)/4}\norm{f}_{Y_\alpha}.
\end{equation}
In particular, no fourth derivative needs to be made small for this
commutator estimate.

\begin{proposition}[Uniform compact right inverse]\label{glue:inverse}
There are $\beta_0,\eta_0,C_D>0$ such that, for
$0<\beta\le\beta_0$ and $0<T\le\eta_0\beta^4$, the operator
$\Pcal$ has a real invariant right inverse
$D:Y_\alpha\to X_\alpha$ with $\norm D\le C_D$.
No initial value is prescribed for $Df$.
\end{proposition}

\begin{proof}
Define
\[
 D_{\mathrm{app}}f=V_i+V_o,\qquad \mathcal R=R_i+J_o.
\]
Equations~\eqref{glue:inner-identity}, \eqref{glue:outer-source}, and
\eqref{glue:outer-identity} give the exact operator identity
\begin{equation}\label{glue:parametrix}
 \Pcal D_{\mathrm{app}}=I+\mathcal R,\qquad
 \norm{D_{\mathrm{app}}}_{Y_\alpha\to X_\alpha}\le C,\qquad
 \norm{\mathcal R}_{Y_\alpha\to Y_\alpha}
 \le C\bigl(\beta^2+\eta^{(1-\alpha)/4}\bigr).
\end{equation}
Choose $\beta_0$ and then $\eta_0$ small enough that the last bound is
at most $1/2$ throughout the asserted range, decreasing the earlier
$\eta_0$ if necessary.  The convergent Neumann series defines
\begin{equation}\label{glue:neumann}
 D=D_{\mathrm{app}}(I+\mathcal R)^{-1}.
\end{equation}
It satisfies $\Pcal D=I$, is real and invariant, and has uniform norm
at most twice the bound for $D_{\mathrm{app}}$.  This is a right inverse
construction and requires no assertion of injectivity.
\end{proof}

\subsection{Smallness of the corrected error}

The right inverse need not have small norm.  What is needed for the
nonlinear problem is its small value on the particular error $E_c$ in
the weaker space $X'_\alpha$.

\begin{lemma}\label{glue:special-error}
For the right inverses of Proposition~\ref{glue:inverse},
\begin{equation}\label{glue:DE-small}
 \norm{DE_c}_{X'_\alpha}
 \le C\left(\beta^{2-2\delta}
       +\eta^{(\alpha'-\alpha)/4}
       +\eta^{(1-\alpha)/4}\right).
\end{equation}
In particular this norm becomes arbitrarily small by choosing $\beta$
and then $\eta=T/\beta^4$ sufficiently small.
\end{lemma}

\begin{proof}
Use the inner output cutoff, rather than the inner input cutoff, to
split
\[
 E_c=f_c+f_e,\qquad f_c=\chi_bE_c,\qquad f_e=(1-\chi_b)E_c.
\]
Equation~\eqref{glue:core-error-small} gives
$\norm{Df_c}_{X_\alpha}\le C\beta^{2-2\delta}$.
The remaining source vanishes for $r\le2\beta$, so
$\chi_i f_e=0$ exactly.  Its inner extension, flat correction,
$R_i f_e$, and $J_i f_e$ all vanish.  Thus
$D_{\mathrm{app}}f_e$ consists solely of the outer zero-data solution
for source $f_e$, followed by multiplication by $\chi_o$, and
$\mathcal R f_e=J_o f_e$.

The extra regularity in Lemma~\ref{glue:background} gives
\begin{equation}\label{glue:outer-error-high}
 \norm{f_e}_{Y_{\alpha'}(s_\beta)}\le C.
\end{equation}
The constant is uniform in $\beta,T$: in the transition annulus all
cutoff derivatives have bounded scaled size, and elsewhere this is
\eqref{glue:error-weight} or the fixed exterior smooth bound.
Apply Lemma~\ref{ana:short} at exponent $\alpha'$ to the outer
zero-data problem with source $f_e$.  The coefficient bounds hold at
that exponent, so its solution has uniformly bounded
$X_{\alpha'}(s_\beta)$ norm.  Equation~\eqref{ana:fourth-small} of
Lemma~\ref{ana:interpolation} gives an $X'_\alpha(s_\beta)$ bound
of $C\eta^{(\alpha'-\alpha)/4}$.  Multiplication by $\chi_o$ is
bounded in these scaled spaces, and $s_\beta$ is comparable to
$\sigma$ on the output support.  This proves the first estimate below;
the commutator bound \eqref{glue:outer-error-small} proves the second:
\begin{equation}\label{glue:app-special-small}
 \norm{D_{\mathrm{app}}f_e}_{X'_\alpha}
 \le C\eta^{(\alpha'-\alpha)/4},\qquad
 \norm{\mathcal Rf_e}_{Y_\alpha}
 \le C\eta^{(1-\alpha)/4}.
\end{equation}
The first estimate includes the temporal H\"older seminorms of all
spatial derivatives through order four, as supplied by
\eqref{ana:fourth-small}.  It makes no claim that the first time
derivative is small: the outer zero-data solution has time derivative
$f_e$ at the initial face.

Finally, \eqref{glue:neumann} gives
$D(I+\mathcal R)=D_{\mathrm{app}}$, so
\begin{equation}\label{glue:special-neumann-identity}
 Df_e=D_{\mathrm{app}}f_e-D\mathcal Rf_e.
\end{equation}
The uniform bound for $D$ and
\eqref{glue:app-special-small} prove the desired estimate for $f_e$.
Combining it with the estimate for $f_c$ proves
\eqref{glue:DE-small}.
\end{proof}

\subsection{The nonlinear correction and finite-time obstruction}

\begin{proof}[Proof of Theorem~\ref{glue:solution}]
Expanding \eqref{glue:equation} at $F$ gives
\begin{equation}\label{glue:nonlinear-equation}
 \Pcal V+E_c+N(V)=0,
\end{equation}
where
\begin{equation}\label{glue:nonlinearity}
 N(V)=v_*^2V^2\A F+
       \bigl(2v_*g_cV+v_*^2V^2\bigr)\A V.
\end{equation}
This expression contains only spatial derivatives of $V$, so is
defined on $X'_\alpha$.  For clarity, on a box of scale $b$, a
function with $\norm{V}_{X'_\alpha}\le\eps$ has scaled sizes
\[
 V=O(\eps b^{-2+2\delta}),\qquad
 \A V=O(\eps b^{-6+2\delta}),
\]
where these statements include the scaled parabolic $C^\alpha$
norms.  By Lemma~\ref{glue:background}, the two quadratic terms in
\eqref{glue:nonlinearity} therefore have size
$C\eps^2b^{-6+4\delta}$, and the cubic term has size
$C\eps^3b^{-6+6\delta}$.  Relative to the $Y$ weight these gain
$b^{2\delta}$ and $b^{4\delta}$, respectively.  Since $b$ is bounded
above on the compact domain, H\"older multiplication and the
factorization of differences of products give uniform bounds
\begin{align}
 \norm{N(V)}_{Y_\alpha}&\le C_N(\eps^2+\eps^3),
 \label{glue:nonlinear-size}\\
 \norm{N(V)-N(W)}_{Y_\alpha}
 &\le C_N(\eps+\eps^2)\norm{V-W}_{X'_\alpha}
 \label{glue:nonlinear-lipschitz}
\end{align}
on the closed ball of radius $\eps$ in $X'_\alpha$.  The fixed sphere
coefficients in $\A$ have bounded rescaled norms, so these estimates
apply on the whole sphere, including the other endpoint $x=1$.

We specify the parameter choices to avoid a dependence on an inverse
norm that is still being chosen.  First fix the profile, exponents,
flat inverse, and the constants $\beta_0,\eta_0,C_D$ of
Proposition~\ref{glue:inverse}.  The nonlinear constant $C_N$ is
uniform throughout this parameter range.  Choose the requested ball
radius $\eps>0$ so small that
\[
 C_D C_N(\eps+\eps^2)\le\tfrac14,\qquad
 C_*\eps\le\tfrac12\mu,
\]
where $C_*$ bounds $\abs{v_*V}/\norm{V}_{X'_\alpha}$ uniformly.
Such a $C_*$ exists because $v_*V$ has scaled size
$C\eps b^{2\delta}$ and $b\le\sqrt2$.
Now take $\beta\le\beta_0$ sufficiently small, and next
$\eta\le\min(\eta_0,a_0^4)$ sufficiently small, that
Lemma~\ref{glue:special-error} gives
$\norm{DE_c}_{X'_\alpha}\le\eps/2$; put $T=\eta\beta^4$.

The map
\[
 \mathcal T(V)=-DE_c-DN(V)
\]
takes the closed $X'_\alpha$ ball of radius $\eps$ into itself by
\eqref{glue:nonlinear-size}, and is a contraction by
\eqref{glue:nonlinear-lipschitz}.  The real invariant subspace is
closed, so the Banach fixed point theorem gives a real invariant
fixed point there.  Its fixed point identity also shows that this
same $V$ belongs to $X_\alpha$, because $E_c+N(V)\in Y_\alpha$ and
$D:Y_\alpha\to X_\alpha$.  Applying $\Pcal D=I$ therefore yields
\eqref{glue:nonlinear-equation} in the classical H\"older sense and
proves \eqref{glue:equation}.  This recovers the time derivative after
the contraction has been performed.

By \eqref{glue:background-positive} and the radius choice,
\[
 \tfrac12\mu\le 1+v_*(F+V)\le M+\tfrac12\mu,
\]
which proves \eqref{glue:positive}.  On each compact time subinterval
of $(-T,0)$, the weight is bounded above and below, and the equation
for $H$ is a uniformly parabolic scalar equation with principal
multiplier $(1+v_*H)^2$.  Its current regularity is
$C^{4+\alpha,1+\alpha/4}$.  Interior spatial difference quotients in
this equation have the same principal operator; their sources are
bounded in $C^\alpha$ because the coefficient is a smooth function
of $x,H$ and $H$ already has four spatial derivatives.  The interior
estimate of Lemma~\ref{ana:schauder}, followed by passage to the
difference-quotient limit on a smaller cylinder, gains one spatial
derivative.  Repetition gains every spatial derivative.  The equation
then gives corresponding spatial bounds for $H_t$; differentiating
the equation in time and repeating the interior estimates gives all
mixed derivatives.  Thus $H$ is smooth on $\Sigma\times(-T,0)$.
Invariant smoothness gives smoothness in $x\in[0,1]$ by
Lemma~\ref{geo:operators}.  At $x=0$, the $X'$ weight reads
\[
 \abs{V(t,0)}\le\eps\sigma(t,0)^{-2+2\delta}
                =\eps\rho^{-1+\delta},
\]
up to the fixed equivalent-norm constant.  This proves
\eqref{glue:center} and finishes the construction.
\end{proof}

\begin{proof}[Proof of Theorem~\ref{thm:main}]
Apply Theorem~\ref{glue:solution} to the nonconstant profile of
Theorem~\ref{prof:profile}.  Proposition~\ref{geo:reduction} turns
the smooth positive coefficient $H$ into a Calabi flow in the
Fubini--Study class on $\CP^{10}$.  Choose its slice at $t=-T/2$ as
the initial metric and translate time so that this slice has time
zero.  It is a smooth K\"ahler metric in a class containing the
Fubini--Study metric.  At the distinguished point,
\[
 H(t,0)=\rho^{-1}G'(0)+O(\rho^{-1+\delta}),\qquad G'(0)<0.
\]
Lemma~\ref{geo:obstruction} shows that the maximal smooth existence
time is exactly $T_*=T/2$ and that, in the elapsed time $\tau$,
\[
 S(\omega(\tau))(p)\sim a(T_*-\tau)^{-1/2},\qquad
 a=-110G'(0)>0.
\]
The metric is $U(10)$ invariant by the radial reconstruction.
These are all the assertions of Theorem~\ref{thm:main}.
\end{proof}

\begin{proof}[Proof of Corollary~\ref{cor:product-singularities}]
For a product metric the determinant factors and the inverse is block
diagonal, so scalar curvatures add.  The product volume form also shows
that their averages add.  Writing \(c_\eta=S(\eta)\), we obtain
\[
 S(\Omega(t))=\pi_X^*S(\omega(t))+c_\eta,
 \qquad \bar S_\Omega=\bar S_\omega+c_\eta.
\]
Thus \(\pi_X^*\phi(t)\), initially zero, satisfies
\eqref{eq:calabi-flow} for \(\Omega_0\).  Smooth local
uniqueness~\cite[Theorem~3.2]{ChenHe2008} identifies this explicit
solution with the unrestricted flow on \([0,T_*)\).  The displayed
scalar-curvature identity gives the stated asymptotic.  A smooth
extension across \(T_*\) would have bounded scalar curvature on the
compact product, so the maximal time is exactly \(T_*\).

For the final assertion put \(m=n-10\).  When \(m>0\), choose
\(Y=\CP^m\) and \(\eta=\frac{m+1}{11}\omega_{\mathrm{FS},m}\).
In our normalization,
\(\operatorname{Ric}(\omega_{\mathrm{FS},j})
=(j+1)\omega_{\mathrm{FS},j}\).  Constant scaling leaves the Ricci
form unchanged, so \(\omega_{\mathrm{FS},10}\) and \(\eta\) both
have Einstein constant \(11\).  Their product form is therefore
K\"ahler--Einstein and lies in \([\Omega_0]\).
For \(m=0\), use Theorem~\ref{thm:main} itself, with the second
factor a point.
\end{proof}
